% Generated draft: requires mathematical and source review.
\documentclass[11pt]{article}
\usepackage[margin=25mm]{geometry}
\usepackage{fontspec}
\setmainfont{Latin Modern Roman}
\usepackage{amsmath,amssymb,mathtools,unicode-math}
\setmathfont{Latin Modern Math}
\usepackage{enumitem,needspace,xurl,hyperref}
\setlength{\parindent}{0pt}
\setlength{\parskip}{5pt}
\setlength{\emergencystretch}{3em}
\newcommand{\N}{\mathbb{N}}
\newcommand{\Z}{\mathbb{Z}}
\newcommand{\Q}{\mathbb{Q}}
\newcommand{\R}{\mathbb{R}}
\newcommand{\C}{\mathbb{C}}
\newcommand{\F}{\mathbb{F}}
\newcommand{\A}{\mathbb{A}}
\newcommand{\PP}{\mathbb{P}}
\newcommand{\E}{\mathbb{E}}
\newcommand{\eps}{\varepsilon}
\newcommand{\dd}{\,\mathrm{d}}
\begin{document}
\section*{UnsolvedMath ID: 10400038}
\textbf{Source problem code:} catalogue entry
\subsection{Definitions and mathematical statement}
Let \(I\) be the interval skeleton. On the apparent reading of the extracted notation, \(A(I)^{(d)}\) is the source's degree-\(d\) chord-diagram module and \(\overrightarrow A(I)^{(d)}\) its arrow-diagram counterpart. Arrow diagrams orient chords, but their exact relation ideal and coefficient ring must be specified separately. Write \(FI\) and \(\overrightarrow{FI}\) for the source's framing-independence submodules, pending their omitted definitions.

The two maps in Polyak's conjecture are schematically
\[
\Phi_d:A(I)^{(d)}\longrightarrow\overrightarrow A(I)^{(d)},
\qquad
\overline\Phi_d:A(I)^{(d)}/FI\longrightarrow
\overrightarrow A(I)^{(d)}/\overrightarrow{FI}.
\]
The assertion is injectivity of both, for every intended degree:
\[
\ker\Phi_d=0,\qquad\ker\overline\Phi_d=0.
\]
Equivalently, once the quotient map is well-defined,
\[
\Phi_d(x)=0\Longrightarrow x=0,
\qquad
\Phi_d(x)\in\overrightarrow{FI}\Longrightarrow x\in FI.
\]
Well-definedness of the second map requires the source rule to carry \(FI\) into \(\overrightarrow{FI}\); this is not an added injectivity conclusion.

The local defining rule was pictorial and is corrupted to an arrow and plus sign in the extraction. It cannot be uniquely reconstructed, including orientations, signs, coefficients, and skeleton relations. In particular, a possible orientation-summing rule is not silently chosen. The coefficient ring and whether \((d)\) denotes exactly a homogeneous degree rather than a truncation are also missing. The draft therefore preserves both precise kernel targets conditional on the actual source maps, without inventing those maps. The supplied research summary merely calls the conjecture not obviously resolved, not a verified modern status theorem.
\subsection{Short English statement}
Prove injectivity of Polyak's two diagram-to-arrow-diagram maps, before and after framing independence. The extracted defining picture, coefficient ring, relations, and grading convention are missing, so the maps cannot yet be fully reconstructed.
\subsection{Sources}
\begin{itemize}
\item Local UnsolvedMath registry, permanent ID \(10400038\), exact source code \texttt{catalogue entry}, original Conjecture 2.16, attributed to M. Polyak.
\item Ohtsuki, \emph{Problems on invariants of knots and 3-manifolds} (2002), Conjecture 2.16, recorded PDF page 39; \url{https://msp.org/gtm/2002/04/gtm-2002-04-024s.pdf}. Recorded access 2026-07-29; pdf-font-structure extraction.
\item The source's pictorial map rule and definitions of the arrow-module/framing relations are not present in recoverable form. The OPEN-TRIAGE research summary supplies no additional precise theorem or publication citation.
\end{itemize}
\end{document}
